% ============================================================
\section{Conclusion}
\label{sec:conclusion}
% ============================================================

We set out to measure how much fine-tuning damages a time-series foundation
model's pre-trained capability, and found that on the standard benchmarks the
capability is mostly not there to damage.  Screened against a lookback-96
linear regression fit on the target data, 5 of 32 cells show the released
checkpoint beating that baseline, all five on one dataset; on three of those
five, unconstrained fine-tuning improves MSE by 31--42\%; and no cell
meets a definition of degradation we fixed in advance.  The practical order of
operations follows directly: before spending effort on LoRA, EWC, L2-SP or a
frozen encoder, check that the features beat a linear fit on your own data, and
if they do, settle the question by running against a frozen encoder and scoring
both on windows you did not select on.

What survives as a positive finding is the dissociation: across the 31
intervention cells the outcome spans $-51.8\%$ to $+46.1\%$ against zero-shot
with no reliable signature in the drift reading (within-backbone Spearman
$\rho{=}{+}0.168$, clustered CI $[-0.469,{+}0.717]$).  How far the encoder moved
does not tell you whether moving it helped.  We had hoped to add a cheap rule
for telling the regimes apart in advance; what we can report instead is that
the premise the rule was built on does not hold on the benchmarks the field
uses to test it.
